\scanpage{8}

\mspage{426}

p h y s i q u e : alors même que je me sentais souvent fatigué et sans entrain avant de me mettre au travail - plus ce travail était ``dur'', me faisant manier une lourde pioche ou de grosses pierres disons, plus après coup je me sentais en forme, lourd d'une bonne fatigue. Et il y avait aussi le contact avec des c h o s e s\quad v i v a n t e s : les plantes qu'il fallait soigner ; la terre qu'il fallait préparer pour les accueillir, puis pailler ou biner ; les aliments qu'il fallait préparer et que je mangeais avec autant de plaisir que j'avais eu à préparer le repas ; le chat réclamant sa pitance, et sa part d'affection ; les ustensiles et outils divers aussi, et jusqu'aux cailloux raboteux et souvent mal léchés qu'il fallait tourner et retourner dans tous les sens, afin de les assembler en des murets qui veuillent bien tenir debout...

Effort physique et contact des choses vivantes - ce sont là deux aspects justement qui manquent dans le travail intellectuel, et qui font qu'un tel travail est par nature incomplet, parcellaire, et à la limite, s'il n'est complété et compensé par autre chose, dangereux voire néfaste. C'est la troisième fois, en à peine plus que trois ans, que j'ai eu l'occasion de m'en rendre compte. Il est même devenu bien clair maintenant, que je suis placé devant une échéance draconienne : changer un certain mode de vie, retrouver un équilibre où le pôle yin de mon être, mon corps, ne soit constamment négligé au profit du pôle yang, l'esprit ou (pour mieux dire) la tête - ou sinon, y laisser ma peau dans les toutes prochaines années. C'est ce que mon corps m'a dit, aussi clairement qu'il est possible de le dire ! J'en arrive à présent à un point dans ma vie où le besoin d'une certaine ``sagesse'' élémentaire est devenu une question de s u r v i e, au sens propre et littéral du terme. C'est sûrement une bonne chose- autrement la dite ``sagesse'' se voyait perpétuellement renvoyée aux calendes, au profit de cette espèce de boulimie dans l'activité intellectuelle, qui a été une des forces dominantes dans toute ma vie d'adulte.

Placé devant une échéance aussi claire : ``changer ou crever !'' - je n'ai pas eu à me sonder pour connaître mon choix. C'est ce qui fait que pendant près de quatre mois, j'aie pu, sans jamais avoir l'impression de me faire violence, m'abstenir de toute activité intellectuelle, maths ou pas maths. J'ai su, sans avoir eu à me le dire, qu'à la limite, un jardinier vivant vaut encore mieux qu'un mathématicien mort (ou un ``philosophe'' ou un ``écrivain'' mort, qu'à cela ne tienne !). Avec un peu de malice, on pourrait ajouter : et même mieux qu'un mathématicien vivant ! (Mais ça, c'est encore une autre histoire...).

Je ne crois pas d'ailleurs que je me trouve acculé un jour à une telle

% archive sheet number: 7

